\section{Introduction}\label{sec:intro}

Suppose you, the player, are shown a grid of points with $h$ rows and $m$ columns. Each
point is either white or blacked out, and a blacked-out point cannot tell you anything.
Before the game starts you choose which points to black out. We call the set of
blacked-out points the \emph{blackout}. Then the game master draws a rectangle with its
four corners on grid points. You cannot see the rectangle. You are only told which of
its corners are white. The rectangle does not have to line up with the grid. Tilted
rectangles count too. You must be able to say where all four corners are with complete
confidence, whatever rectangle the game master draws. Kagey's question is how many
points you can black out and still do this.

The order of play matters. The blackout is fixed before the rectangle is drawn, so a
blackout is only good if no two rectangles could ever look the same to you.

\begin{figure}[t]
\centering
\resizebox{0.95\textwidth}{!}{\begin{tikzpicture}[x=0.9cm,y=0.9cm,baseline=(current bounding box.north)]
  \begin{scope}[xshift=0.000cm]
    \draw[black,line width=0.36pt,dash pattern=on 1.35pt off 1.35pt] (1,-1) -- (4,-1) (1,-2) -- (4,-2) (1,-3) -- (4,-3) (1,-1) -- (1,-3) (2,-1) -- (2,-3) (3,-1) -- (3,-3) (4,-1) -- (4,-3);
    \filldraw[fill=white,draw=black,line width=0.54pt] (1,-1) circle (0.3);
    \filldraw[fill=white,draw=black,line width=0.54pt] (2,-1) circle (0.3);
    \filldraw[fill=white,draw=black,line width=0.54pt] (3,-1) circle (0.3);
    \filldraw[fill=white,draw=black,line width=0.54pt] (4,-1) circle (0.3);
    \filldraw[fill=white,draw=black,line width=0.54pt] (1,-2) circle (0.3);
    \filldraw[fill=white,draw=black,line width=0.54pt] (2,-2) circle (0.3);
    \filldraw[fill=kyellow,draw=black,line width=0.54pt] (3,-2) circle (0.3);
    \filldraw[fill=white,draw=black,line width=0.54pt] (4,-2) circle (0.3);
    \filldraw[fill=white,draw=black,line width=0.54pt] (1,-3) circle (0.3);
    \filldraw[fill=white,draw=black,line width=0.54pt] (2,-3) circle (0.3);
    \filldraw[fill=white,draw=black,line width=0.54pt] (3,-3) circle (0.3);
    \filldraw[fill=white,draw=black,line width=0.54pt] (4,-3) circle (0.3);
    \node[lab] at (1,-0.45) {1};
    \node[lab] at (2,-0.45) {2};
    \node[lab] at (3,-0.45) {3};
    \node[lab] at (4,-0.45) {4};
    \node[lab] at (0.45,-1) {1};
    \node[lab] at (0.45,-2) {2};
    \node[lab] at (0.45,-3) {3};
    \node[capt,anchor=north] at (2.50,-3.45) {(a) the point $(3,2)$};
  \end{scope}
  \begin{scope}[xshift=4.680cm]
    \draw[black,line width=0.36pt,dash pattern=on 1.35pt off 1.35pt] (1,-1) -- (4,-1) (1,-2) -- (4,-2) (1,-3) -- (4,-3) (1,-1) -- (1,-3) (2,-1) -- (2,-3) (3,-1) -- (3,-3) (4,-1) -- (4,-3);
    \filldraw[fill=kgray,draw=black,line width=0.54pt] (1,-1) circle (0.3);
    \draw[black,line width=0.54pt] (0.788,-0.788) -- (1.212,-1.212) (0.788,-1.212) -- (1.212,-0.788);
    \filldraw[fill=kgray,draw=black,line width=0.54pt] (2,-1) circle (0.3);
    \draw[black,line width=0.54pt] (1.788,-0.788) -- (2.212,-1.212) (1.788,-1.212) -- (2.212,-0.788);
    \filldraw[fill=kgray,draw=black,line width=0.54pt] (3,-1) circle (0.3);
    \draw[black,line width=0.54pt] (2.788,-0.788) -- (3.212,-1.212) (2.788,-1.212) -- (3.212,-0.788);
    \filldraw[fill=kgray,draw=black,line width=0.54pt] (4,-1) circle (0.3);
    \draw[black,line width=0.54pt] (3.788,-0.788) -- (4.212,-1.212) (3.788,-1.212) -- (4.212,-0.788);
    \filldraw[fill=kgray,draw=black,line width=0.54pt] (1,-2) circle (0.3);
    \draw[black,line width=0.54pt] (0.788,-1.788) -- (1.212,-2.212) (0.788,-2.212) -- (1.212,-1.788);
    \filldraw[fill=white,draw=black,line width=0.54pt] (2,-2) circle (0.3);
    \filldraw[fill=white,draw=black,line width=0.54pt] (3,-2) circle (0.3);
    \filldraw[fill=white,draw=black,line width=0.54pt] (4,-2) circle (0.3);
    \filldraw[fill=kgray,draw=black,line width=0.54pt] (1,-3) circle (0.3);
    \draw[black,line width=0.54pt] (0.788,-2.788) -- (1.212,-3.212) (0.788,-3.212) -- (1.212,-2.788);
    \filldraw[fill=white,draw=black,line width=0.54pt] (2,-3) circle (0.3);
    \filldraw[fill=white,draw=black,line width=0.54pt] (3,-3) circle (0.3);
    \filldraw[fill=white,draw=black,line width=0.54pt] (4,-3) circle (0.3);
    \node[lab] at (1,-0.45) {1};
    \node[lab] at (2,-0.45) {2};
    \node[lab] at (3,-0.45) {3};
    \node[lab] at (4,-0.45) {4};
    \node[lab] at (0.45,-1) {1};
    \node[lab] at (0.45,-2) {2};
    \node[lab] at (0.45,-3) {3};
    \node[capt,anchor=north] at (2.50,-3.45) {(b) our blackout};
  \end{scope}
  \begin{scope}[xshift=9.360cm]
    \draw[black,line width=0.36pt,dash pattern=on 1.35pt off 1.35pt] (1,-1) -- (4,-1) (1,-2) -- (4,-2) (1,-3) -- (4,-3) (1,-1) -- (1,-3) (2,-1) -- (2,-3) (3,-1) -- (3,-3) (4,-1) -- (4,-3);
    \draw[kred,line width=2.55pt,line join=round] (2,-2) -- (3,-1) -- (4,-2) -- (3,-3) -- cycle;
    \filldraw[fill=kgray,draw=black,line width=0.54pt] (1,-1) circle (0.3);
    \draw[black,line width=0.54pt] (0.788,-0.788) -- (1.212,-1.212) (0.788,-1.212) -- (1.212,-0.788);
    \filldraw[fill=kgray,draw=black,line width=0.54pt] (2,-1) circle (0.3);
    \draw[black,line width=0.54pt] (1.788,-0.788) -- (2.212,-1.212) (1.788,-1.212) -- (2.212,-0.788);
    \filldraw[fill=kgray,draw=black,line width=0.54pt] (3,-1) circle (0.3);
    \draw[black,line width=0.54pt] (2.788,-0.788) -- (3.212,-1.212) (2.788,-1.212) -- (3.212,-0.788);
    \filldraw[fill=kgray,draw=black,line width=0.54pt] (4,-1) circle (0.3);
    \draw[black,line width=0.54pt] (3.788,-0.788) -- (4.212,-1.212) (3.788,-1.212) -- (4.212,-0.788);
    \filldraw[fill=kgray,draw=black,line width=0.54pt] (1,-2) circle (0.3);
    \draw[black,line width=0.54pt] (0.788,-1.788) -- (1.212,-2.212) (0.788,-2.212) -- (1.212,-1.788);
    \filldraw[fill=white,draw=black,line width=0.54pt] (2,-2) circle (0.3);
    \filldraw[fill=white,draw=black,line width=0.54pt] (3,-2) circle (0.3);
    \filldraw[fill=white,draw=black,line width=0.54pt] (4,-2) circle (0.3);
    \filldraw[fill=kgray,draw=black,line width=0.54pt] (1,-3) circle (0.3);
    \draw[black,line width=0.54pt] (0.788,-2.788) -- (1.212,-3.212) (0.788,-3.212) -- (1.212,-2.788);
    \filldraw[fill=white,draw=black,line width=0.54pt] (2,-3) circle (0.3);
    \filldraw[fill=white,draw=black,line width=0.54pt] (3,-3) circle (0.3);
    \filldraw[fill=white,draw=black,line width=0.54pt] (4,-3) circle (0.3);
    \node[lab] at (1,-0.45) {1};
    \node[lab] at (2,-0.45) {2};
    \node[lab] at (3,-0.45) {3};
    \node[lab] at (4,-0.45) {4};
    \node[lab] at (0.45,-1) {1};
    \node[lab] at (0.45,-2) {2};
    \node[lab] at (0.45,-3) {3};
    \node[capt,anchor=north] at (2.50,-3.45) {(c) the rectangle};
  \end{scope}
  \begin{scope}[xshift=14.040cm]
    \draw[black,line width=0.36pt,dash pattern=on 1.35pt off 1.35pt] (1,-1) -- (4,-1) (1,-2) -- (4,-2) (1,-3) -- (4,-3) (1,-1) -- (1,-3) (2,-1) -- (2,-3) (3,-1) -- (3,-3) (4,-1) -- (4,-3);
    \filldraw[fill=kyellow,draw=black,line width=0.54pt] (2,-2) circle (0.3);
    \filldraw[fill=kyellow,draw=black,line width=0.54pt] (4,-2) circle (0.3);
    \filldraw[fill=kyellow,draw=black,line width=0.54pt] (3,-3) circle (0.3);
    \node[lab] at (1,-0.45) {1};
    \node[lab] at (2,-0.45) {2};
    \node[lab] at (3,-0.45) {3};
    \node[lab] at (4,-0.45) {4};
    \node[lab] at (0.45,-1) {1};
    \node[lab] at (0.45,-2) {2};
    \node[lab] at (0.45,-3) {3};
    \node[capt,anchor=north] at (2.50,-3.45) {(d) what we are shown};
  \end{scope}
\end{tikzpicture}
}
\caption{One round on the $3\times4$ grid. Black dots are blacked out and open dots are
white. In (d) the grey specks are points we are told nothing about.}
\label{fig:game}
\end{figure}

Figure~\ref{fig:game} plays one round on the $3\times4$ grid. We number the columns
from the left and the rows from the top, and we write $(i,j)$ for the point in column
$i$ and row $j$, as in Figure~\ref{fig:game}(a). In (b) we black out all of row $1$
and all of column $1$, which is six points. In (c) the game master draws the tilted
square with corners $(2,2)$, $(3,1)$, $(4,2)$ and $(3,3)$. The corner $(3,1)$ is in row
$1$, so it is black, and in (d) we are shown only $(2,2)$, $(4,2)$ and $(3,3)$. That is
enough. Three corners of a rectangle always determine the fourth
(Lemma~\ref{lem:threecorners}), and here the fourth corner has to be $(3,1)$.
Section~\ref{sec:three} shows that this blackout works against every rectangle on the
$3\times4$ grid, and that six is the most anyone can black out there.

Kagey poses the problem on the web page of his Open Problem
Collection~\cite{kagey001}, and it is Problem 1 of the printed collection he handed
out at the fourteenth Gathering 4 Gardner~\cite{kageyG4G14}. The page lists seven
variants of the question and no results. We write $f(h,m)$ for the largest size of a
blackout of the $h\times m$ grid that keeps every rectangle identifiable, and $a_h(m)$
for the number of blackouts of that size. Blackouts are counted with their
coordinates, so two blackouts that differ by a symmetry of the grid count twice. These
two quantities are the tables \seqnum{A397315} and \seqnum{A399580} in the
OEIS~\cite{oeis}, read by antidiagonals with both indices starting at $2$. The entries
were created in September 2026 from the computations of Section~\ref{sec:verify}. The
comments on \seqnum{A397315} record the strip values $f(2,m)=m+1$, the formula
$f(3,m)=m+2$ as a conjecture checked up to width $10$, and the guess that
$f(h,m)=h+m$ whenever $\min(h,m)\ge4$, apart from $f(4,4)=7$ and $f(4,5)=8$. Both
statements are settled in this paper. The first is a theorem, with a proof that
covers every number of rows, and the second is false.

The three-row formula was proved first elsewhere. A machine-checked proof of exactly
this statement, with the same rectangles and the same rules, was merged into the
\texttt{trureturing} repository of The Omega Institute on 15 September
2026~\cite{trureturing2026}, before the argument of this paper was written. That
proof counts the black points of the three rows by inclusion and exclusion. The proof
here is the case $h=3$ of a bound for every number of rows
(Theorem~\ref{thm:rowbound}), and we do not claim the three-row value as new.

The smallest case already shows what matters. On two rows every rectangle is
axis-aligned, because a tilted rectangle has corners in at least three rows
(Section~\ref{sec:strip}). So the $2\times3$ strip has three rectangles, one for each
pair of columns, and we write $R_{a,b}$ for the rectangle on columns $a<b$.

\begin{example}[the $2\times3$ strip]\label{ex:strip23}
Take the blackout with the four black points $(3,1)$, $(1,2)$, $(2,2)$ and $(3,2)$,
which keeps $(1,1)$ and $(2,1)$ white. The rectangle $R_{1,2}$ shows $(1,1)$ and
$(2,1)$, the rectangle $R_{1,3}$ shows $(1,1)$ alone, and $R_{2,3}$ shows $(2,1)$
alone. The three views differ, so this blackout is good. Now black out columns $1$ and
$2$ instead, keeping $(3,1)$ and $(3,2)$. Then $R_{1,3}$ and $R_{2,3}$ both show
exactly the two kept points. If we are shown $(3,1)$ and $(3,2)$, we cannot tell
whether the left side is in column $1$ or in column $2$. Both blackouts have four
points, so the size alone decides nothing. Can we black out five points? Then only
one point is kept, and every view is either empty or that one point. That is two
possible views for three rectangles, so two of them agree. So the maximum blackout of
the $2\times3$ strip is $4$, that is, $f(2,3)=4$. Counting the good blackouts of
size four gives $a_2(3)=12$ (Theorem~\ref{thm:strip}).
\end{example}

The example contains the two mechanisms of the whole paper. Two rectangles look the
same exactly when every corner that belongs to only one of them is black
(Lemma~\ref{lem:diff}), and a rectangle that shows three corners is never confused
with another one (Lemma~\ref{lem:threecorners}). So everything turns on the
rectangles that show at most two corners.

\paragraph{Results.} The paper proves the following six statements. Each is restated
with its numbered proof in the body, and the finite computer checks are marked where
they occur.

\begin{thmA}[narrow grids]
For all $h,m\ge1$ we have $f(h,m)\ge h+m-1$, and for $m\ge3$ we have
$f(h,m)\le m+\lfloor h^2/4\rfloor$. For $m\ge3$ the values on two and three rows
are $f(2,m)=m+1$ and $f(3,m)=m+2$. On four rows, $f(4,m)=m+4$ for every $m\ge6$,
while $f(4,4)=7$ and $f(4,5)=8$. On five rows, $f(5,5)=10$ and $f(5,m)=m+6$ for every
$m\ge21$. [[P2-fiverow]]
\end{thmA}

The lower bound is the boundary cross of Lemma~\ref{lem:cross} and the upper bound is
the row bound of Theorem~\ref{thm:rowbound}. The strip value is
Theorem~\ref{thm:strip}, the three-row value Theorem~\ref{thm:threevalue}, the
four-row values Theorem~\ref{thm:fourvalue} and Corollary~\ref{cor:fourboundary}, the
$5\times5$ value Theorem~\ref{thm:fivefive}, the five-row family
Theorem~\ref{thm:fiverow} and the small five-row widths
Proposition~\ref{prop:fivesmall}. The values $f(4,4)$ and $f(4,5)$ rest on a finite
computer check of $72$ and $1{,}440$ candidates. So the excess $f(h,m)-m$ over the
width is $1$, $2$ and $4$ on two, three and four rows, and it is $6$ on five rows once
the width is at least $21$. The guess $f(h,m)=h+m$ predicts excess $5$ on five rows,
and it fails from width $21$ on (Theorem~\ref{thm:refute}).

\begin{thmB}[three-row optima]
For $m\ge4$,
\begin{align*}
a_3(m)&=4(m-1)(3m+2)\,3^{m-5}+E_m,\\
0&\le E_m\le4\binom m2\kappa^3\rho^{m-2}+m\kappa^2\rho^{m-1},
\end{align*}
where $\rho=1+\sqrt3$, $\kappa=9/\rho^2$, and $E_m$, defined in
Section~\ref{sec:three}, counts the optima in which rows $1$ and $3$ are black
together in some column. [[P2-Em]] The generating function $\sum_{m\ge3}a_3(m)z^m$ is
rational, and its reduced denominator is exactly
$(1-3z)^3(1-2z-2z^2)^3(1+2z^2)^3$. Hence $a_3(m)\sim\frac4{81}m^23^m$, and the
fraction of the $m^23^{m-1}$ candidates that pass the axis-aligned test and are valid
tends to $4/27$.
\end{thmB}

The formula is Theorem~\ref{thm:threecount}, the generating function
Theorem~\ref{thm:threegf} with Appendix~\ref{app:threegf}, and the asymptotics
Corollary~\ref{cor:threeasymp}. The sequence $a_3(m)$ starts
$53,224,892,3420,12704$ at $m=3$, which is row $3$ of \seqnum{A399580}, and it
satisfies a linear recurrence of order $15$ from $m=19$ on
(Corollary~\ref{cor:threerec}). In words, almost every three-row optimum has one
black point in each column plus two more points, one in a column black in rows $1$
and $2$ and one in a column black in rows $2$ and $3$, and the constant $4/27$ is one
third, the share of these optima among the candidates, times $(6/9)^2$, the share of
the nine ways to fill the two columns beside each extra point that survive the tilted
squares there.

\begin{thmC}[four rows]
Let $m\ge4$. No valid blackout of the $4\times m$ grid hides all four corners of a
rectangle. For $m\ge6$,
\[
a_4(m)\le3\,(m)_4\,4^{m-4}\Bigl(\frac{53}{64}\Bigr)^{\max(0,\lfloor m/4\rfloor-4)},
\qquad (m)_4=m(m-1)(m-2)(m-3),
\]
so $a_4(m)=O(m^4\,212^{m/4})$, and $\sum_{m\ge6}a_4(m)z^m$ is a rational function.
The first values are $a_4(6)=26$ and $a_4(7)=292$, found by computer. [[P2-fourrow]]
\end{thmC}

The hidden-rectangle statement is Theorem~\ref{thm:hidden}, the bound
Theorem~\ref{thm:foursurvival}, the rationality Theorem~\ref{thm:fourgf} and the
values Proposition~\ref{prop:fourvalues}. The bound says nothing below $m=20$. Its
point is the base $212^{1/4}<4$, which says that the fraction of the $3(m)_44^{m-4}$
candidates passing the axis-aligned test that survive the tilted rectangles tends to
zero exponentially, unlike the three-row fraction $4/27$. By
Corollary~\ref{cor:rarity} the same happens at every height $h\ge4$ for the
blackouts that attain the row bound.

\begin{thmD}[finite collar]
Let $h\ge3$, put $W_h=1+\lfloor(h-1)^2/4\rfloor$, and let $S$ be a blackout of an
$h\times m$ grid with $m\ge1$. Choose the top row or the bottom row as the padding
row, and choose one side of the grid or both. For $L,R\ge0$ let $S(L,R)$ be the
blackout of the $h\times(L+m+R)$ grid obtained by adding $L$ columns on the left and
$R$ columns on the right, each black exactly in the padding row, with $L=0$ or $R=0$
on a side that was not chosen. Then $S(L,R)$ is valid for all such $L$ and $R$ if and
only if it is valid when $L$ and $R$ equal $W_h$ on each chosen side. In particular,
the blackout of the $5\times21$ grid in Figure~\ref{fig:witness21} stays valid when
padded on the left by columns black in row $1$, so $f(5,m)=m+6$ for every $m\ge21$,
and the conjecture $f(h,m)=h+m$ for $\min(h,m)\ge5$ is false.
\end{thmD}

The collar statement is Theorem~\ref{thm:collar}, the five-row family is
Theorem~\ref{thm:fiverow} with its case list in Appendix~\ref{app:fivecases}, and the
refutation is Theorem~\ref{thm:refute}. The padding row has to be a boundary row
(Remark~\ref{rem:interior}), and the same witness cannot be padded on the right by
more than three columns (Remark~\ref{rem:rightpad}). The threshold $21$ is sufficient,
and we do not claim it is the smallest width with excess six. In words, a tilted
rectangle that shows at most two corners has a black corner above the padding row, so
it lives within $W_h$ columns of the original grid, and all that a far-away padding
column can do is complete an axis-aligned rival, which the same pair of rows already
completes inside the test grid.

\begin{thmE}[every fixed height]
Fix $h\ge1$ and encode a blackout of the $h\times m$ grid as the word
$B_1B_2\cdots B_m$ of its column supports, whose letters are subsets of
$\{1,\dots,h\}$. The valid words form a regular language, and the series
$\sum_{m\ge0}\sum_S z^mu^{|S|}$, summed over the valid blackouts $S$ of width $m$, is
a rational function with constant term $1$. If $h\ge2$, then for $m\ge3$ the excess
$f(h,m)-m$ lies between $h-1$ and $\lfloor h^2/4\rfloor$, it is eventually periodic
in $m$, and $\sum_{m\ge0}a_h(m)z^m$ is a rational function, with $a_h(0)=1$.
\end{thmE}

Regularity is Theorem~\ref{thm:regular}, the bivariate series
Theorem~\ref{thm:rational}, and the periodicity and the optimum series
Theorem~\ref{thm:periodic}. Validity is not a property of bounded windows of columns
(Proposition~\ref{prop:nonlocal}), so the recognizer has to carry pairs of rectangles
rather than local patterns. Theorem~E explains the shape of Theorems~B and~C, but it
computes neither the period nor the threshold nor the excess, and it does not say
whether the row bound is attained on six or more rows.

\begin{thmF}[squares]
There are constants $c>0$ and $n_0$ such that
$f(n,n)\ge c\,n^{6/5}(1+\ln n)^{-2/5}$ for all $n\ge n_0$, and
$f(n,n)\le2n^{4/3}+2n$ for all $n\ge3$.
\end{thmF}

This is Theorem~\ref{thm:square}, proved in Appendix~\ref{app:square}. The threshold
$n_0$ is not effective, so the theorem produces no explicit square on which
$f(n,n)-2n$ is large. The upper bound is elementary, and it is weaker than Neuwirth's
bound for bipartite graphs of girth eight (Remark~\ref{rem:neuwirth}). The two
regimes differ. For fixed $h$ the excess $f(h,m)-m$ stays below $\lfloor h^2/4\rfloor$,
while on squares $f(n,n)-2n$ is unbounded, and Theorem~E cannot be specialized to
$h=m$ because its thresholds depend on $h$.

\paragraph{Prior work and what is new.} The reduction of the game to a hitting set
(Corollary~\ref{cor:hitting}) is the test cover formulation.
Katona~\cite{katona1966} studies separating systems of a finite set,
Bellitto~\cite[Section~2.1]{bellitto2018} writes the separating code problem as a
hitting set on symmetric differences, and Foucaud~\cite[Section~1.1]{foucaud2015}
distinguishes test covers, in which one empty signature is allowed, from
discriminating codes, in which none is. Our difference lemma is this reduction for
rectangles. On the graph side, a blackout whose row-column graph has no cycle of
length four or six keeps the axis-aligned rectangles apart (Lemma~\ref{lem:girth}),
so the axis-aligned part of the problem asks for the largest bipartite graph with
given parts and girth at least eight. Erd\H{o}s~\cite{erdos1080} asked about this
function, de Caen and Sz\'ekely~\cite{decaenszekely1997} bound it for dense graphs,
and Neuwirth~\cite[Theorem~1]{neuwirth2001} proves $e\le v+\lfloor w^2/4\rfloor$ for
parts of sizes $v$ and $w$ under the hypothesis $v\ge\lfloor w^2/4\rfloor$. Our row
bound (Theorem~\ref{thm:rowbound}) has the same value $m+\lfloor h^2/4\rfloor$ for
every $m\ge3$, with a proof by one tree in each column and Mantel's
theorem~\cite{mantel1907}, and its equality structure (Theorem~\ref{thm:equality})
is what the counting results use. Naor and Verstra\"ete~\cite{naorverstraete2005} and
Gy\H{o}ri~\cite{gyori1997} bound bipartite graphs without six-cycles alone. None of
these sources sees the tilted rectangles, and the tilted rectangles are what makes
the counts in Theorems~B and~C smaller than the graph counts. Counting the words
accepted by a finite automaton by a rational generating function is classical
(Flajolet and Sedgewick~\cite[Propositions~I.3 and~V.6]{flajolet2009}), and column
transfer is the usual tool for fixed-height grid problems in this journal, for
independent sets of grid graphs (Euler~\cite{euler2005}), for juggling patterns
(Butler, Choi, Kim, and Seo~\cite{butler2019}) and for domination polynomials of
grids (Mertens~\cite{mertens2024}), as well as for domination numbers of grid graphs
(Alanko, Crevals, Isopoussu, \"Osterg\r{a}rd, and Pettersson~\cite{alanko2011}).
Jiang~\cite{jiang2018} proves that optimal identifying codes in strips can be taken
periodic. Guilhot, Little, and Parkinson~\cite[Theorem~1]{glp2025} prove that the
words of a regular language attaining the maximum of a bounded additive weight form
a regular language, and the eventual periodicity of the powers of a max-plus matrix
is the cyclicity theorem of Baccelli, Cohen, Olsder, and Quadrat~\cite{baccelli1992}.
What these results do not supply is the recognizer, because validity here is not
local (Proposition~\ref{prop:nonlocal}). What is new in this paper is the exact
count of the three-row optima with its generating function and its constant, the
hidden-rectangle theorem, the forbidden blocks and the finite memory on four rows,
the finite collar theorem with the five-row family, and the paired-rectangle
recognizer behind Theorem~E. The square bounds use the alteration method of Alon and
Spencer~\cite[Section~3.2]{alonspencer2016}, Spencer's bound for independent sets in
uniform hypergraphs~\cite{spencer1972} in the form stated by Conlon, Fox, and
Sudakov~\cite[Section~2]{conlonfoxsudakov2016}, and the parametrization of lattice
rectangles by a primitive direction, which Babichev and
Babichev~\cite{babichev2026mfcs}, Babichev and Shpakova~\cite{babichevshpakova2026}
and Babichev and Pinchuk~\cite{babichevpinchuk2026} use to count all rectangles in a
grid.

\paragraph{Organization.} Section~\ref{sec:conv} fixes the conventions and proves the
difference lemma and the three-corner lemma. Section~\ref{sec:strip} treats strips.
Section~\ref{sec:graph} builds the row-column graph, proves the row bound with its
equality structure and the boundary cross, and lists the facts about tilted
rectangles that every later section uses. Sections~\ref{sec:three}, \ref{sec:four}
and~\ref{sec:five} treat three, four and five rows, and Section~\ref{sec:height}
treats every fixed height and then the squares. Section~\ref{sec:verify} lists what
was computed, what is formalized in Lean, and where the code and data
are~\cite{repo}. Section~\ref{sec:open} collects the open questions.
Appendix~\ref{app:fivecases} has the case list for the five-row family,
Appendix~\ref{app:threegf} the remaining three-row generating functions, and
Appendix~\ref{app:square} the proofs for squares.
